% =============================================================================
% exercises/ejercicios_cap09.tex
% Ejercicios propuestos — Capítulo 9: Estratosfera
% =============================================================================

\section*{Ejercicios Propuestos}
\addcontentsline{toc}{section}{Ejercicios propuestos}
\label{sec:ejercicios_cap09}

\begin{bloque_ejercicios}[Ejercicios --- Cap\'itulo 9: Estratosfera]

\begin{ejercicios}


% ----- Ejercicio 5: Reacciones de halógenos y competencia RH1 vs RH2 ------
\item Los \'atomos de hal\'ogeno siguen dos rutas competitivas
(reacciones \ref{RH1} y \ref{RH2}):
\[
  \ce{X + RH -> R^. + HX} \quad\text{(RH1)} \qquad
  \ce{X + O3 -> XO^. + O2} \quad\text{(RH2)}
\]

\begin{enumerate}[label=\alph*)]
    \item Seg\'un el cap\'itulo, las fracciones que reaccionan
          v\'ia \ref{RH2} son: F~$\sim$0\%, Cl~50\%,
          Br~99\%, I~100\%. Ordene los hal\'ogenos de mayor
          a menor reactividad frente a compuestos con H
          y explique la tendencia en t\'erminos de energ\'ia
          de enlace H--X.

    \item El \'atomo de Cl liberado de \ce{CH3Cl} sigue la
          secuencia:
          \ce{CH3Cl + OH^. -> ^.CH2Cl + H2O},
          \ce{^.CH2Cl + O2 ->[M] CH2ClOO^.},
          \ce{CH2ClOO^. + NO -> NO2 + HCHO + Cl^.}.
          Identifique en esta secuencia: el paso que regenera
          Cl, el que convierte NO a \ce{NO2} y el producto
          carbonilo formado.

    \item Explique por qu\'e los CFCs, a pesar de ser m\'as
          pesados que el aire, alcanzan la estratosfera.
          ?`Qu\'e propiedad fisicoqu\'imica les permite
          sobrevivir el tr\'ansito por la troposfera intactos?

    \item La reacci\'on \ce{N2O5 + NaX_{(s)} -> XNO2 + NaNO3_{(s)}}
          ocurre en aerosoles de sal marina. Explique por qu\'e
          esta reacci\'on \textbf{heterog\'enea} es m\'as
          r\'apida que la correspondiente en fase gas, y
          qu\'e implicaciones tiene para la liberaci\'on
          de hal\'ogenos activos en la atm\'osfera marina.
\end{enumerate}

% ----- Ejercicio 6: Ciclos de Chapman y balance de ozono estratosférico ----
\item Las reacciones de Chapman (\ref{OE1}--\ref{OE3}) describen
el balance natural del ozono estratosf\'erico.

\begin{enumerate}[label=\alph*)]
    \item Escriba las cuatro reacciones de Chapman
          (\ref{OE1}, \ref{OE2}, \ref{OE3} y \ce{O3 + O -> 2O2})
          e identifique cu\'ales producen y cu\'ales destruyen ozono.

    \item Sume las reacciones de producci\'on neta y las de
          destrucci\'on neta del ciclo completo y verifique
          que el resultado neto global es
          \ce{2O2 ->[$h\nu$] O2 + O + O3 -> 2O2}
          (ciclo nulo de ox\'igeno).

    \item Aplique la aproximaci\'on de estado seudoestacionario
          (PSSA) al \'atomo O para demostrar que:
          \[
              [\ce{O}] = \frac{k_{\ref{OE1}}[\ce{O2}]}{
              k_{\ref{OE2}}[\ce{O2}][\mathrm{M}] +
              k_4[\ce{O3}]}
          \]
          donde $k_4$ es la constante de la reacci\'on
          \ce{O3 + O -> 2O2}.

    \item En la estratosfera real, las concentraciones de ozono
          medidas son \textbf{menores} que las predichas solo
          por Chapman. Mencione los tres ciclos catal\'iticos
          del cap\'itulo que explican esta discrepancia
          e indique la familia radical responsable de cada uno.
\end{enumerate}
% ----- Ejercicio 7: Balance de ozono estratosférico (Chapman) --------------
\item El balance natural de ozono en la estratosfera se describe mediante
las reacciones de Chapman.

\begin{enumerate}[label=\alph*)]
    \item Escriba las cuatro reacciones del mecanismo de Chapman:
          dos de formaci\'on (reacciones \ref{OE1} y \ref{OE2})
          y dos de destrucci\'on (reacci\'on \ref{OE3} y su reacci\'on
          inversa de regeneraci\'on). Indique las longitudes de onda
          de los fotones involucrados en cada fotolisis.

    \item Usando la analog\'{\i}a del ``balde que gotea'' descrita en
          el cap\'{\i}tulo, explique qu\'e condici\'on debe cumplirse para
          que el nivel de ozono permanezca constante, y qu\'e ocurre
          cuando se introducen compuestos adicionales como los CFC.

    \item El cap\'{\i}tulo indica que la concentraci\'on promedio de ozono
          es 300~DU, equivalente a una capa de \SI{3}{\milli\metre}.
          Si en el ``agujero'' ant\'artico cae a 100~DU:
          \begin{enumerate}[label=\roman*)]
              \item ?`Cu\'antos mil\'{\i}metros de espesor
                    representan 100~DU?
              \item Usando la definici\'on de que 1~DU $\equiv$
                    $2.69\times10^{16}$~mol\'eculas~cm$^{-2}$,
                    calcule el n\'umero total de mol\'eculas de
                    \ce{O3} en una columna de 1~cm$^2$ de base
                    para 300~DU y para 100~DU.
          \end{enumerate}

    \item La columna total de aire comprimido a STP tiene un espesor
          de \SI{8}{\kilo\metre}. Calcule qu\'e porcentaje del espesor
          total de la atm\'osfera representa la capa de ozono de 300~DU
          (\SI{3}{\milli\metre}).
\end{enumerate}

\textit{[Respuesta parcial:
(c-i) $100\,\text{DU} \rightarrow \SI{1}{\milli\metre}$;
(c-ii) $300\,\text{DU}: N = 300 \times 2.69\times10^{16}
      = 8.07\times10^{18}$~mol\'eculas~cm$^{-2}$;
      $100\,\text{DU}: N = 2.69\times10^{18}$~mol\'eculas~cm$^{-2}$;
(d) $3\,\text{mm}/8\,000\,000\,\text{mm} = 3.75\times10^{-7}$, es
    decir, $\sim 0.000038\%$ del espesor total]}

% ----- Ejercicio 8: Ciclos catalíticos de pérdida de ozono -----------------
\item Los ciclos catal\'{\i}ticos \ce{HO_$x$}, \ce{NO_$x$} y
\ce{ClO_$x$} son responsables de la p\'erdida acelerada de ozono
en la estratosfera.

\begin{enumerate}[label=\alph*)]
    \item Para el ciclo \ce{HO_$x$}, escriba las dos reacciones
          propagadoras (\ce{OH + O3} y \ce{HO2 + O3}) y la reacci\'on
          global neta. ?`Cu\'al es el cat\'alizador del ciclo?
          ?`C\'omo se termina el ciclo?

    \item Para el ciclo catal\'{\i}tico del \ce{NO_$x$}
          (reacciones \ref{R10}--\ref{R11}), escriba las dos
          reacciones y la reacci\'on neta. Calcule cu\'antas
          mol\'eculas de \ce{O_$x$} (equivalente a \ce{O3})
          se consumen por cada ciclo completo.

    \item El ciclo \ce{ClO_$x$} se inicia con la fotolisis del CFC-12
          (\ce{CF2Cl2}). Escriba:
          (i)~la reacci\'on de fotolisis del CFC-12,
          (ii)~las dos reacciones propagadoras del ciclo
          \ce{Cl}/\ce{ClO}, y
          (iii)~la reacci\'on neta global.

    \item Compare los tres ciclos catal\'{\i}ticos (\ce{HO_$x$},
          \ce{NO_$x$}, \ce{ClO_$x$}) llenando la siguiente tabla:

{%
\centering
\begin{tabular}{lccc}
\toprule
\textbf{Caracter\'{\i}stica} & \ce{HO_$x$} & \ce{NO_$x$} & \ce{ClO_$x$} \\
\midrule
Radical cat. & & & \\
Fuente principal & & & \\
Reacci\'on de terminaci\'on & & & \\
Reservorio inactivo & & & \\
\bottomrule
\end{tabular}
\par
}%
\end{enumerate}
% ----- Ejercicio 7: Ciclos catalíticos de pérdida de ozono ----------------
\item Los ciclos de p\'erdida catal\'itica del ozono en la
estratosfera involucran las familias \ce{HO_$x$},
\ce{NO_$x$} y \ce{ClO_$x$}.

\begin{enumerate}[label=\alph*)]
    \item Complete la siguiente tabla de ciclos catal\'iticos
          usando los mecanismos del \textbf{Cuadro~\ref{desozo}}
          y las secciones \ref{roh}--\ref{rclo}:

{%
\centering
\begin{tabular}{llll}
\toprule
\textbf{Ciclo} & \textbf{Paso 1} & \textbf{Paso 2} & \textbf{Neto} \\
\midrule
\ce{HO_$x$}   & \ce{OH + O3 -> HO2 + O2}  & \underline{\hspace{2cm}} & \ce{2O3 -> 3O2} \\
\ce{NO_$x$}   & \underline{\hspace{2.5cm}} & \ce{NO2 + O^. -> NO + O2} & \ce{O3 + O^. -> 2O2} \\
\ce{ClO_$x$}  & \ce{Cl + O3 -> ClO + O2}  & \underline{\hspace{2cm}} & \ce{O3 + O^. -> 2O2} \\
\bottomrule
\end{tabular}
\par
}%

    \item El ciclo~2 del \textbf{Cuadro~\ref{desozo}} involucra
          la dimeraci\'on \ce{ClO + ClO -> (ClO)2}. ?`Por qu\'e
          este ciclo es especialmente importante en el
          \textit{agujero ant\'artico} y no en latitudes medias?
          Relacione con las condiciones de temperatura y
          concentraci\'on de Cl.

    \item El ciclo~3 involucra tanto Cl como Br:
          \ce{ClO + BrO -> Cl + Br + O2}.
          Explique por qu\'e la presencia simult\'anea de
          cloro y bromo amplifica sinerg\'isticamente
          la destrucci\'on del ozono.

    \item Un solo \'atomo de Cl puede destruir hasta
          $10^5$ mol\'eculas de \ce{O3}. Si la fotolisis de
          1~kg de CFC-12 (\ce{CF2Cl2}, masa molar = 121~g/mol)
          libera todos sus \'atomos de Cl disponibles,
          calcule el n\'umero de mol\'eculas de \ce{O3}
          que podr\'ian destruirse.
          ($N_A = 6.022\times10^{23}$~mol\'ec/mol)
\end{enumerate}

\textit{[Respuesta:
(d) moles CFC-12 = $1000/121 = 8.26$~mol;
    \'atomos Cl = $2\times8.26\times6.022\times10^{23}
    = 9.95\times10^{24}$~\'atomos;
    mol\'ec \ce{O3} destruidas $= 9.95\times10^{24}
    \times 10^5 = 9.95\times10^{29}$~mol\'ec]}

% ----- Ejercicio 8: Reservorios de NOx y ClOx en la estratosfera ----------
\item Las familias \ce{NO_$y$} y \ce{ClO_$y$} incluyen radicales
activos y especies reservorio no radicales.

\begin{enumerate}[label=\alph*)]
    \item Construya un diagrama de interconversiones para
          la familia \ce{NO_$y$}, incluyendo:
          \ce{NO}, \ce{NO2} (familia \ce{NO_$x$}),
          \ce{HNO3}, \ce{N2O5} (reservorios diurnos y nocturnos)
          y las reacciones que los interconvierten seg\'un
          el cap\'itulo.

    \item Explique por qu\'e el \ce{HNO3} activa al \ce{NO_$x$}
          en la \textbf{estratosfera} (mediante fot\'olisis y
          reacci\'on con OH), pero en la \textbf{troposfera}
          es un sumidero terminal (eliminado por lluvia).
          ?`Qu\'e propiedad determina esta diferencia?

    \item Para la familia \ce{ClO_$y$}, distinga entre:
          (i) \ce{ClO_$x$} (radicales activos: Cl y ClO),
          (ii) reservorios HCl y \ce{ClNO3}.
          Calcule el tiempo de vida de HCl en la estratosfera
          usando $k_{\ce{HCl+OH}} = 8\times10^{-13}$
          cm$^3$~mol\'ec$^{-1}$~s$^{-1}$ y
          $[\ce{OH}] = 5\times10^5$~mol\'ec/cm$^3$.

    \item El \ce{N2O} es transportado a la estratosfera
          sin sumideros troposf\'ericos. Solo el 5\% de su
          p\'erdida estratosf\'erica produce \ce{NO_$x$}
          activo. Si las emisiones globales de \ce{N2O}
          son $\approx 18$~Tg~N/a\~no, estime la producci\'on
          anual de \ce{NO} activo en la estratosfera.
\end{enumerate}

\textit{[Respuesta:
(c) $\tau_{\ce{HCl}} = 1/(8\times10^{-13}
    \times 5\times10^5)
    = 2.5\times10^6$~s $\approx$ 29~d\'ias;
(d) $0.05\times18 = 0.9$~Tg~N/a\~no]}

% ----- Ejercicio 9: Unidades Dobson y columna de ozono --------------------
\item La Unidad Dobson (DU) mide la columna total de ozono:
1~DU $\equiv$ $2.69\times10^{16}$~mol\'ec/cm$^2$.

\begin{enumerate}[label=\alph*)]
    \item La columna normal de ozono es 300~DU.
          Calcule el n\'umero total de mol\'eculas de
          \ce{O3} en una columna de 1~cm$^2$ de base.

    \item Convierta 300~DU a espesor f\'isico equivalente
          (en mm) a \SI{0}{\celsius} y 1~atm, sabiendo
          que 1~DU $= 0.01$~mm de ozono puro comprimido.

    \item Durante el agujero ant\'artico, la columna cae
          a $\approx$~100~DU. Calcule:
          (i) el n\'umero de mol\'eculas de \ce{O3}
          perdidas por cm$^2$ respecto a la columna normal,
          (ii) la masa de \ce{O3} perdida por cada
          \si{\kilo\metre\squared} de superficie ant\'artica
          afectada (masa molar \ce{O3} = 48~g/mol).

    \item Relacione el adelgazamiento de la capa de ozono
          con el incremento de radiaci\'on UV-B que llega
          a la superficie: por cada 1\% de disminuci\'on
          en \ce{O3}, la UV-B aumenta $\approx$~2\%.
          ?`Cu\'anto aumenta la UV-B cuando la columna
          pasa de 300 a 100~DU?
\end{enumerate}

\textit{[Respuesta:
(a) $300 \times 2.69\times10^{16}
    = 8.07\times10^{18}$~mol\'ec/cm$^2$;
(b) $300 \times 0.01 = 3$~mm;
(c-i) $(300-100)\times2.69\times10^{16}
    = 5.38\times10^{18}$~mol\'ec/cm$^2$;
(d) reducci\'on = $(300-100)/300 = 66.7\%$;
    aumento UV-B $= 66.7 \times 2 = 133\%$]}
% ----- Ejercicio 9: PSCs y agujero de ozono --------------------------------
\item Las nubes estratosf\'ericas polares (PSC) desempe\~nan un papel
fundamental en el agujero de ozono ant\'artico.

\begin{enumerate}[label=\alph*)]
    \item El cap\'{\i}tulo indica que las PSC se forman alrededor de
          \SI{-72}{\celsius}. Convierta esta temperatura a Kelvin
          y explique por qu\'e las PSC son m\'as frecuentes sobre
          la Ant\'artida que sobre el \'Artico.

    \item En la superficie de las PSC los gases de ``dep\'osito''
          \ce{HCl} y \ce{ClONO2} se activan. Escriba la reacci\'on
          heterog\'enea de \ce{N2O5} con \ce{NaCl} en aerosol salino
          que produce una especie activa de cloro (\ce{XNO2}) y una sal
          s\'olida, seg\'un el cap\'{\i}tulo.

    \item La reacci\'on \ce{N2O + O(^1D) -> 2NO} en la estratosfera
          es responsable del 5\% de la p\'erdida del \ce{N2O}.
          Si la tasa total de p\'erdida de \ce{N2O} es
          $F = \SI{1.0e10}{\kilo\gram}/$a\~no, ?`cu\'anta masa
          (en kg/a\~no) se convierte a NO por esta v\'{\i}a?
          ?`Qu\'e ocurre con el 95\% restante?

    \item Molina y Rowland recibieron el Premio Nobel de Qu\'{\i}mica
          en 1995 junto con Paul Crutzen. Explique la contribuci\'on
          espec\'{\i}fica de cada uno en la comprensi\'on de la
          destrucci\'on del ozono estratosf\'erico, con base en la
          informaci\'on del cap\'{\i}tulo.
\end{enumerate}

\textit{[Respuesta parcial:
(a) \SI{-72}{\celsius} $= \SI{201}{\kelvin}$;
(c) $5\%$ de $10^{10}$ kg/a\~no $= 5\times10^8$~kg/a\~no;
    el 95\% se convierte en \ce{N2} (inerte) por fotolisis y  oxidaci\'on con \ce{O(^1D)}]}

% ----- Ejercicio 10 (avanzado): Ciclo 2 de ClO y ecuación de continuidad ---
\item \textbf{[Avanzado]} El cap\'{\i}tulo describe tres ciclos de
destrucci\'on de ozono con hal\'ogenos (\textbf{Cuadro~\ref{desozo}}).

\begin{enumerate}[label=\alph*)]
    \item Para el \textbf{Ciclo 2} (\ce{ClO + ClO -> (ClO)2 -> \ldots}),
          escriba los cuatro pasos elementales y la reacci\'on global
          neta. Identifique cu\'al paso requiere un fot\'on ($h\nu$)
          y cu\'al produce el \'atomo de Cl libre necesario para
          continuar el ciclo.

    \item Para el \textbf{Ciclo 3} (que involucra tanto Cl como Br),
          escriba los cuatro pasos y la reacci\'on neta.
          Compare la eficiencia del Ciclo 3 con la del Ciclo 1
          en t\'erminos de la raz\'on de mol\'eculas de \ce{O3}
          destruidas por \'atomo de hal\'ogeno.

    \item La ecuaci\'on de continuidad del ozono estratosf\'erico es:
          \[
              \frac{\partial[\ce{O3}]}{\partial t}
              = P - L - \nabla\cdot(\vec{v}[\ce{O3}])
          \]
          Identifique cada t\'ermino f\'{\i}sicamente e indique
          qu\'e proceso representa cada uno:
          \begin{enumerate}[label=\roman*)]
              \item el t\'ermino $P$,
              \item el t\'ermino $L$,
              \item el t\'ermino $\nabla\cdot(\vec{v}[\ce{O3}])$.
          \end{enumerate}
          ?`En qu\'e regi\'on de la Tierra domina el transporte
          advectivo sobre la producci\'on qu\'{\i}mica?

    \item La circulaci\'on de Brewer-Dobson transporta \ce{O3} desde
          los tr\'opicos hacia los polos. Explique por qu\'e, a pesar
          de que el \ce{O3} se produce principalmente en la estratosfera
          tropical, las m\'aximas concentraciones se observan en
          latitudes altas (salvo en la Ant\'artida en primavera austral).
\end{enumerate}


% ----- Ejercicio 10 (avanzado): Integración de ciclos y PSC ---------------
\item \textbf{[Avanzado]} El agujero ant\'artico de ozono
se forma por la combinaci\'on de los tres ciclos del
\textbf{Cuadro~\ref{desozo}} activados en la superficie
de las nubes estratosf\'ericas polares (PSC).

\begin{enumerate}[label=\alph*)]
    \item Las PSC se forman cuando la temperatura
          estratosf\'erica cae por debajo de
          $\approx$ \SI{-72}{\celsius}. Explique por qu\'e
          el hemisferio sur forma PSC con mayor facilidad
          que el norte, y en qu\'e \'epoca del a\~no
          ocurre el m\'aximo agotamiento de ozono.

    \item En la superficie de las PSC, los reservorios
          de cloro inactivos se convierten en formas reactivas:
          \[
            \ce{HCl + ClNO3 ->[\text{PSC}] Cl2 + HNO3}
          \]
          Explique por qu\'e esta reacci\'on heterog\'enea
          es cr\'itica: ?`qu\'e le ocurre al \ce{Cl2} cuando
          regresa la luz solar en primavera? ?`Y al \ce{HNO3}?

    \item La ecuaci\'on de continuidad del ozono
          estratosf\'erico es:
          \[
            \frac{\partial [\ce{O3}]}{\partial t}
            = P - L - \nabla\cdot(\vec{v}[\ce{O3}])
          \]
          En condiciones de agujero ant\'artico,
          $P \approx 0$ (oscuridad invernal) y el
          transporte $\nabla\cdot(\vec{v}[\ce{O3}]) < 0$
          (\'ozono entra al vort\'ex polar).
          ?`Cu\'al de los ciclos 1, 2 o 3 domina el
          t\'ermino $L$ durante la primavera austral?
          Justifique con las concentraciones t\'ipicas
          de Cl y la temperatura.

    \item El Protocolo de Montreal (1987) redujo la
          producci\'on de CFCs. Si la concentraci\'on
          estratosf\'erica de Cl alcanz\'o su m\'aximo
          de $\approx$~3.5~ppb en los a\~nos 2000 y
          decrece a raz\'on de $\approx$~0.7\%/a\~no,
          estime en qu\'e a\~no la concentraci\'on
          volver\'ia a los niveles preindustriales
          ($\approx$~0.6~ppb). Comente si es realista
          dado el tiempo de vida de los CFCs.
\end{enumerate}

\textit{[Respuesta parcial:
(d) reducci\'on necesaria: $(3.5-0.6)/3.5 = 82.9\%$;
    a 0.7\%/a\~no: $\ln(3.5/0.6)/0.007
    \approx 254$~a\~nos desde el m\'aximo $\approx$ a\~no 2254;
    coherente con vidas de CFCs de 50--1700~a\~nos
    (Cuadro cap\'itulo 5)]}

\end{ejercicios}

\end{bloque_ejercicios}
% =============================================================================
% FIN DE ejercicios_cap09.tex
% =============================================================================
